% CV de recherche. Dates, statuts et contributions mis à jour par l'auteur le 5 novembre 2026.
% Mise en page du CV de recherche approuvé; titres bibliographiques publiés avec l'autorisation de l'auteur.
\documentclass[11pt,letterpaper]{article}
\usepackage[utf8]{inputenc}
\usepackage[T1]{fontenc}
\usepackage[lf]{ebgaramond}
\usepackage[margin=0.68in,top=0.6in,bottom=0.62in]{geometry}
\usepackage{microtype}
\usepackage{enumitem}
\usepackage{titlesec}
\usepackage{fancyhdr}
\usepackage{xcolor}
\usepackage[hidelinks,pdfauthor={Shayan Hajhashemi},pdftitle={CV de recherche — Shayan Hajhashemi}]{hyperref}
\setlength{\parindent}{0pt}
\setlength{\parskip}{2pt}
\setlist[itemize]{leftmargin=1.15em,itemsep=1pt,topsep=2pt,parsep=0pt}
\titleformat{\section}{\large\bfseries}{}{0em}{}[\vspace{-4pt}\titlerule]
\titlespacing*{\section}{0pt}{9pt}{4pt}
\pagestyle{fancy}
\fancyhf{}
\renewcommand{\headrulewidth}{0pt}
\fancyfoot[L]{\scriptsize Shayan Hajhashemi · Mise à jour : 5 novembre 2026}
\fancyfoot[R]{\scriptsize\thepage}
\setlength{\footskip}{20pt}
\newlength{\projectheadingwidth}
\newcommand{\project}[2]{%
  \vspace{2pt}%
  \settowidth{\projectheadingwidth}{\textbf{#1}\quad\textit{#2}}%
  \ifdim\projectheadingwidth>\linewidth
    \textbf{#1}\par\nopagebreak[4]\textit{#2}\par\nopagebreak[4]%
  \else
    \textbf{#1}\hfill\textit{#2}\par\nopagebreak[4]%
  \fi}
\newcommand{\me}{\textbf{Hajhashemi, S.}}
\begin{document}
\begin{center}
{\LARGE\bfseries Shayan Hajhashemi}\\[3pt]
Modèles du monde · Apprentissage de représentations · Systèmes dynamiques et commande\\[3pt]
Université McGill · Montréal, Canada\quad
\href{https://mila.quebec/en/directory/shayan-hajhashemi}{Profil Mila}\\[2pt]
\small\href{https://shayshay42.github.io/}{Site web}\quad
\href{https://github.com/shayshay42}{GitHub}\quad
\href{https://scholar.google.com/citations?user=uW3aYOcAAAAJ}{Google Scholar}\quad
\href{https://ca.linkedin.com/in/hshay}{LinkedIn}\\[2pt]
\footnotesize shayan [dot] hajhashemi [at] mail [dot] mcgill [dot] ca
\end{center}
\vspace{-11pt}

\section*{Profil de recherche}
Chercheur en apprentissage automatique scientifique : apprentissage de représentations et systèmes dynamiques non linéaires. Mes travaux couvrent les modèles du monde multimodaux, les solveurs inverses appris et la commande optimale sous contraintes en simulation différentiable. Intérêt pour les représentations conditionnées par l'action et les modèles du monde fiables pour l'apprentissage robotique et la planification.
Recherche de stages et de postes de recherche à temps plein.

\section*{Formation}
\textbf{Université McGill} — Doctorat, sciences quantitatives de la vie\hfill Prévu en 2026\\
Direction : Morgan Craig et Amin Emad. Outils d'apprentissage automatique scientifique pour la dynamique non linéaire et la commande, appliqués à la pharmacologie.\\
\textbf{Université McGill} — B.Sc. (Honours), immunologie interdépartementale\hfill 2021\\
Projet Honours : motifs de Turing dans les systèmes de réaction--diffusion; direction : Anmar Khadra.

\section*{Travaux de recherche sélectionnés}
\project{Modèles du monde multimodaux : alignement, prédiction et commande}{Soumis à AISTATS 2027}
\begin{itemize}
\item Co-conçu des approches d'alignement des représentations et de commande prédictive non linéaire (NMPC) pour les modèles du monde multimodaux; contribué à l'analyse théorique.
\item Comparé notre régularisation des représentations à SINDy et aux modèles de fondation de séries temporelles (TSFM), sur la même tâche de NMPC du système de Lorenz.
\end{itemize}

\project{WeLD : inférence d'EDP par distillation en formulation faible sans étiquettes}{Manuscrit de travail}
\begin{itemize}
\item Coécrit une méthode sans étiquettes qui apprend des poids sur les équations faibles : elle distille des estimations classiques sur de longues fenêtres d'observation pour inférer les paramètres d'EDP sur de courtes fenêtres.
\item Analysé le biais du modèle enseignant et les corrections par variables instrumentales; exploré des modèles du monde latents qui apprennent l'inférence des paramètres par prévision, sans enseignant.
\end{itemize}

\project{WAR : identification de systèmes par formulation faible et régularisation adaptative}{Soumis à AISTATS 2027}
\begin{itemize}
\item Associé sélection apprise de fenêtres temporelles et régularisation adaptative à des solveurs différentiables sous contraintes en formulation faible; estimé des paramètres non linéaires à partir d'observations bruitées par projection variable.
\item Démontré de meilleures performances en inférence inverse sur des systèmes de test issus de six familles d'EDO, avec évaluations distinctes des paramètres et des trajectoires. Testé la sensibilité au bruit et à l'échantillonnage peu dense ou irrégulier.
\end{itemize}

\project{OIL : apprentissage d'initialisations par l'optimisation en aval}{Soumis à ICLR 2027}
\begin{itemize}
\item Entraîné des initialiseurs par passage de messages à travers cinq étapes différentiables d'ajustement de trajectoires, avec une perte à trois termes sur les paramètres initiaux, les paramètres affinés et les trajectoires simulées.
\item Analysé la géométrie des sensibilités par l'information de Fisher, les échecs de simulation et les diagnostics a posteriori.
\end{itemize}

\project{DiffDose : commande optimale différentiable pour la médecine personnalisée}{Prépublication; en évaluation à npj Systems Biology and Applications}
\begin{itemize}
\item Optimisé doses bornées et instants d'administration par simulation différentiable d'EDO et d'équations à retard. Comparé différentiation automatique, adjoints analytiques et différences finies en commande en boucle ouverte.
\end{itemize}

\project{Dynamique latente et apprentissage de représentations}{Prépublication; en évaluation à Nature Computational Science}
\begin{itemize}
\item Coécrit un cadre variationnel contraignant les représentations apprises d'états cellulaires (RNA-seq unicellulaire) à des flots sur un paysage énergétique, avec des EDO neuronales à gradient de potentiel pour prédire trajectoires et destins cellulaires.
\end{itemize}

\newpage
\section*{Publications et manuscrits}
\begin{enumerate}[leftmargin=1.2em,itemsep=7pt,topsep=2pt,parsep=0pt]
\item \me; Shum-Tim, C.; Craig, M.; Emad, A.
\textit{WAR: Weak-form Adaptive Regularization for Dynamical Systems Inverse Problems.}
Soumis le 7 octobre 2026. En cours d'évaluation à \textit{AISTATS 2027}.
\item \me; Shum-Tim, C.; Prakash, S.; Emad, A.
\textit{Predators, Prey, and Sloppy Parameters: Fine-Tuning Amortized Inference for Generalized Lotka--Volterra.}
Soumis le 18 septembre 2026. En cours d'évaluation à \textit{ICLR 2027}. \textbf{Co-premier auteur.}
\item Ramezanian-Panahi, M.; \me; Rish, I.
\textit{Multimodal World Models: Alignment, Prediction, and Control.}
Soumis le 7 octobre 2026. En cours d'évaluation à \textit{AISTATS 2027}.
\item \me; Emad, A.; Craig, M.
\textit{DiffDose: Differentiable Programming for Personalized Dose-Regimen Optimal Control.}
Soumis le 13 septembre 2026. DOI bioRxiv : \href{https://doi.org/10.64898/2026.09.07.749974}{10.64898/\allowbreak 2026.09.07.749974}.
En cours d'évaluation à \textit{npj Systems Biology and Applications}.
\item Poursina, A.; \me; Mikaeili Namini, A.; Saberi, A.; Emad, A.; Najafabadi, H.~S.
\textit{A Latent Space Thermodynamic Model of Cell Differentiation.}
Soumis le 6 mars 2026. DOI bioRxiv : \href{https://doi.org/10.64898/2026.03.04.709512}{10.64898/\allowbreak 2026.03.04.709512}.
En cours d'évaluation à \textit{Nature Computational Science}.
\item \href{https://www.gene.com/scientists/our-scientists/vincent-lemaire}{Lemaire, V.}; \me; Rotte, A.; Seyed-Ahmadi, A.; Tripathi, S.; Chanu, P.; Wu, B.; Li, C.; Ramanujan, S.; Gadkar, K.; Jin, J.
\textit{Explainability analysis of QSP-derived virtual populations identifies robust pre-treatment biomarkers of immunotherapy.}
Accepté pour publication dans \textit{npj Precision Oncology}.
\end{enumerate}

\section*{Publications complémentaires}
\begin{enumerate}[start=7,leftmargin=1.2em,itemsep=7pt,topsep=2pt,parsep=0pt]
\item Brassard, A.; Dhoparee-Doomah, I.; \me; Bourdeau, F.; Giannias, B.; Spicer, J.; Ferri, L.~E.; Ma, K.; Cools-Lartigue, J.
\textit{Neutrophil Extracellular Trap Deposition Establishes an Immunosuppressive Environment in Non-metastatic Lymph Nodes and Facilitates Heterologous Metastatic Outgrowth.}
28 octobre 2026. bioRxiv, version 1. DOI : \href{https://doi.org/10.64898/2026.10.23.761387}{10.64898/\allowbreak 2026.10.23.761387}.

\item Dhoparee-Doomah, I.; Ma, K.; Brassard, A.; Leo, S.; \me; \ldots; Cools-Lartigue, J.
\textit{ROR$\gamma$ Inverse Agonism Redirects Tumour-Conditioned Emergency Myelopoiesis and Attenuates Established Tumour Progression.}
7 octobre 2026. bioRxiv, version 1. DOI : \href{https://doi.org/10.64898/2026.10.01.755741}{10.64898/\allowbreak 2026.10.01.755741}.
En cours d'évaluation à \textit{Cancer Discovery}.
\item Dhoparee-Doomah, I.; Li, Y.; Al Dow, M.; \me; \ldots; Cools-Lartigue, J.
\textit{A Measurable Systemic Immune Phenotype Links Circulating Myeloid Dysfunction to Tumour Spatial Organization in Gastroesophageal Adenocarcinoma.}
18 septembre 2026. bioRxiv, version 1. DOI : \href{https://doi.org/10.64898/2026.09.17.752366}{10.64898/\allowbreak 2026.09.17.752366}.
\item Roche, P.; \me; Akavia, U.~D. (2020).
\textit{Nanobody administration of Cas9 improves nuclear localization and HDR-enhanced synthetic lethality in A549 lung carcinoma.}
Retiré pour des raisons de propriété intellectuelle. \href{https://www.jenthera.com/}{Jenthera Therapeutics}.
\end{enumerate}
Listes complètes des auteurs et brevet : \href{https://shayshay42.github.io/publications.html}{shayshay42.github.io/publications.html}.

\newpage
\section*{Expérience de recherche}
\textbf{Genentech} — Stagiaire en modélisation et simulation\hfill 2024--2025\\
Direction : Vincent Lemaire; San Francisco, États-Unis.
\begin{itemize}
\item Entraîné des modèles génératifs conditionnels par flow matching et transport optimal pour des populations virtuelles issues de simulateurs; appliqué les distributions apprises à l'analyse de biomarqueurs.
\end{itemize}
\textbf{Expérience antérieure :} Chercheur consultant, Jenthera Therapeutics (2020--2021); associé de recherche, Centre de recherche sur le cancer Goodman (2019--2020). Recherche expérimentale sur les nanocorps et vision par ordinateur.

\section*{Projets techniques sélectionnés}
\textbf{Manipulation robotique en simulation} — \href{https://github.com/shayshay42/robot}{Prototype d'apprentissage robotique en simulation}, 2022--2023
\begin{itemize}
\item Créé des scripts d'entraînement et d'évaluation pour la manipulation simulée avec MuJoCo, robosuite, Gymnasium et Stable-Baselines3; exploré soft actor--critic et le clonage comportemental à partir de démonstrations.
\end{itemize}
\textbf{Apprentissage de représentations de systèmes dynamiques et de graphes} — McGill, 2021--2022
\begin{itemize}
\item Exploré les plongements de séries temporelles par Transformers, autoencodeurs LSTM, caractéristiques de Fourier et ondelettes pour la recherche de nouveauté; étudié les plongements par convolution sur graphes relationnels et la comparaison structurale de motifs d'ARN 3D.
\end{itemize}

\section*{Compétences techniques}
\textbf{Programmation :} Python, Julia, MATLAB; PyTorch, JAX/diffrax, Pyro, Julia SciML.\\
\textbf{Méthodes :} EDO neuronales, réseaux de neurones sur graphes, identification de systèmes, optimisation différentiable, inférence variationnelle, flow matching, analyse de sensibilité, commande prédictive.\\
\textbf{Projets de simulation :} MuJoCo, robosuite, Gymnasium, Stable-Baselines3; expériences d'apprentissage par renforcement et par imitation.

\section*{Présentations}
\begin{itemize}[leftmargin=1.15em,itemsep=6pt]
\item \me\ \href{https://meetings.siam.org/sess/dsp_talk.cfm?p=157552}{\textit{Real-MVP: Virtual Populations As Simulation-Based Conditional Neural Densities.}}
SIAM Conference on the Life Sciences, Cleveland, Ohio, 6 juillet 2026.
Présentation orale invitée à une conférence internationale. MS11: Recent Advances in Digital Twins and Trial Simulations to Guide Clinical Decisions.
\item Ramezanian-Panahi, M.; \me; Rish, I.
\href{https://www.main2026.org/abstracts}{\textit{Cross-modal control in the Virtual Brain.}}
Présentation par affiche (4D), \href{https://www.main2026.org/program}{MAIN 2026 — Montreal AI and Neuroscience}, HEC Montréal, Montréal, Canada, 30 octobre 2026.
\item Shum-Tim, C.; \me\ \textit{Weak-Form Distillation for Label-Free PDE Inverse Problems.}
\href{https://wmw-aspen.github.io/}{World Modeling for Physics Workshop}, Aspen Center for Physics, Aspen, Colorado, 28 février--5 mars 2027.
\textbf{Présentation orale acceptée; à venir.}
\end{itemize}

\section*{Enseignement et service}
\textbf{Enseignement :} Apprentissage automatique scientifique et EDO dans Julia, formateur d'atelier (2023); dynamique non linéaire et méthodes de calcul aux cycles supérieurs, auxiliaire d'enseignement (2022).\\
\textbf{Service :} Évaluateur pour une conférence Julia (2024).

\section*{Prix et distinctions sélectionnés}
BESC D du CRSNG (2023--2026); bourse doctorale FRQ-Santé (attribuée pour 2023--2026); lauréat du concours départemental de thèse en trois minutes (2022); Dean's Multidisciplinary First Class Honours (2019--2021).
\end{document}
